\makeabstract
    {
    Investigating Localized Surface Plasmon Resonance Effects in Hybrid AuNP-TADF Thin Films
    }
    {
    Jonathan R. Mendez\textsuperscript{1}, Christopher A. Nieves\textsuperscript{3}, Batul Shakir\textsuperscript{2}, Aditi Aggarwal\textsuperscript{2}, Jeevesh Magnani\textsuperscript{2}, Bogdan Dragnea\textsuperscript{2}, Sara E. Skrabalak\textsuperscript{2}, Ricardo J. V\'azquez\textsuperscript{2}
    }
    {
    \textsuperscript{1} Department of Chemistry, Amherst College, Amherst, MA\\
\textsuperscript{2} Department of Chemistry, Indiana University Bloomington, Bloomington, IN\\
\textsuperscript{3} Department of Chemistry, University of Puerto Rico at Cayey, Cayey, PR
    }
    {
    Thermally activated delayed fluorescence (TADF) materials exhibit temperature-dependent photophysical changes, as thermal energy facilitates reverse intersystem crossing from the triplet to the singlet excited states, giving rise to delayed fluorescence. Gold nanoparticles (AuNPs) exhibit localized surface plasmon resonance (LSPR), through which absorbed light can generate localized thermal effects. This study investigates how LSPR from AuNPs influences the photophysical behavior of the TADF emitter 2,4,5,6-tetra(9H-carbazol-9-yl)isophthalonitrile (4CzIPN) in polymer thin films, with emphasis on changes in its delayed fluorescence lifetime and its potential as a probe of local temperature changes near AuNPs.
    }
    {
    Polymer thin films containing poly(methyl methacrylate) (PMMA), 4CzIPN, and AuNPs were prepared by spin coating and drop casting and characterized using ultraviolet-visible (UV-Vis) absorption and steady-state emission spectroscopy. Temperature-dependent fluorescence lifetime measurements were used to determine how the lifetime of 4CzIPN changes with temperature. Hybrid AuNP-4CzIPN films were further examined using dark-field microscopy, fluorescence lifetime imaging microscopy (FLIM), and transient absorption spectroscopy to investigate their photophysical response to plasmon excitation.
    }
    {
    The delayed fluorescence lifetime of 4CzIPN in polymer thin films decreased with increasing temperature, establishing a relationship between the photophysical response of 4CzIPN and thermal changes in its environment. In hybrid thin films containing both 4CzIPN and AuNPs, a decrease in fluorescence lifetime was also observed following plasmon excitation, consistent with a possible contribution from LSPR-generated heating to the observed lifetime decrease. However, the observed changes in lifetime cannot yet be attributed exclusively to plasmonic heating, as other interactions between AuNPs and 4CzIPN, including quenching, may also contribute to the measured response.
    }
    {
    These results show the potential for using the temperature-dependent delayed fluorescence lifetime of 4CzIPN to probe local temperature changes in hybrid AuNP-TADF thin films. Further work is needed to separate LSPR-generated heating from other AuNP-4CzIPN interactions, such as quenching, and determine whether TADF emitters can be used as local temperature probes near plasmonic nanoparticles.
    }

    \newpage
    